%% 中文内容审阅稿：版式遵循 CHI 2027 的匿名单栏模板；正式投稿须使用英文 main.tex。
\documentclass[manuscript,review,anonymous]{acmart}
%% Do not load ctex: it changes ACM-controlled line spacing. xeCJK only
%% supplies Chinese glyph selection while retaining acmart's review layout.
\usepackage{xeCJK}
\setCJKmainfont{FandolSong-Regular}
\newcommand{\cn}[1]{#1}
\usepackage{booktabs}
\usepackage{enumitem}
\usepackage{xcolor}
\newcommand{\todo}[1]{\textcolor{red}{[\textbf{待完成：}#1]}}
\setcopyright{none}
\settopmatter{printacmref=false,printccs=false,printfolios=true}

\begin{document}

\title{\cn{谁在使用开发者社区？}\\Agentic Coding \cn{中开发者--Agent 协同轨迹与技术知识生态的角色转变}}

\begin{abstract}
AI Coding 正在改变开发者使用技术社区的方式。开发者不再只是在文档、样例和社区讨论之间自行检索与拼接知识；他们还把检索、代码起草、工具调用和迭代交给 Agent。Agent 表面上的自主性并未消除开发者设定目标、提供本机事实、判断证据和验证结果的责任。本文提出开发者--Agent 协同机制（Developer--Agent Collaboration Mechanism，DACM），将协同开发描述为开发者、Coding Agent、技术知识生态与本地工程环境之间的反馈循环。DACM 以 Agent 的可观察工作方式为基础：上下文建构、检索、候选行动、工具反馈与迭代。研究采用顺序式设计：首先以 26 个 CANN/CUDA 结果目标相近的开发任务开展形成性知识生态审计，以构造具有不同知识供给条件的任务情境；随后计划招募约 20 名开发者，结合近期 AI 协同项目回放与统一算子开发情境任务，观察参与者在自身 IDE、CLI 和常用 Agent 中如何推进任务。研究提出多主体协同轨迹分析，用于追踪证据取得、事实请求、人工接管、验证 checkpoint 与停止点。本文的目标不是比较模型能力或硬件性能，而是解释技术社区如何成为开发者--Agent 协同体的知识基础设施，以及其设计如何影响任务能否凭借可追溯证据推进。
\end{abstract}

\keywords{人机协同；AI Coding；Coding Agent；开发者体验；技术文档；开发者社区}
\maketitle

\section{引言}

当开发者遇到陌生 API、编译错误、版本冲突或性能问题时，传统路径是搜索、阅读文档、比较社区答案、在本机运行，再调整下一步。AI Coding 工具正在重新组织这一过程。开发者可以描述目标或粘贴报错，Agent 随后检索资料、读取工程文件、提出补丁、调用工具、观察反馈并继续修改。这并不使开发者成为被动审批者：开发者依然决定目标与边界、选择提供哪些上下文和权限、评估风险，并对最终交付负责。

因此，变化并不是简单的“AI 做得更多”。Coding Agent 通常运行在一个有边界的循环中：它从任务和可见材料建构上下文，检索或读取证据，提出行动，再在获授权时调用工具，并根据反馈修订。它无法可靠推断未提供或无法取得的事实，例如实际版本组合、硬件拓扑、私有代码语义、业务约束或验收阈值。页面对人可见不代表 Agent 能稳定取得正文；Agent 生成的命令或代码看似流畅，也不代表其前置条件已经在本地环境中成立。

这同时改变了技术社区的角色。文档、样例、兼容性矩阵、仓库、Issue 与论坛不再仅仅是人类阅读的材料；它们成为 Agent 必须发现、访问、解释、引用并转化为行动的证据基础设施。若只评估最终答案质量，就难以区分问题来自模型、检索链路、文档工程、版本治理、社区供给，还是缺失的现场事实。若只评估传统门户体验，也无法解释为何开发者在 Agent 给出建议后仍反复搜索、纠正与验证。

本文提出 DACM，用于研究开发者、Agent、技术知识生态与本地工程环境如何共同推进任务。我们以 Agent 的可观察工作循环作为分析基础，但不把模型内部机制或新的 Agent 架构作为本文主张。本文关注的问题是：\emph{在 Agentic Coding 中，开发者、Agent、技术知识生态和本地工程环境如何共同推进开发任务？}

研究采用顺序式设计。Study 0 在 26 个 CANN/CUDA 开发任务上审计一个具备网页检索与读取能力的 Agent 如何取得官方与第三方知识、版本证据以及可执行步骤；这不是“哪一个生态更强”的结论，而是选择知识供给差异显著的情境。Study 1 回放参与者近期的 AI 协同开发任务。Study 2 则让同一批参与者在自己的 IDE、CLI、浏览器和 Agent 中使用 AddCustom 到 PyTorch 的连续工程情境包。情境包不要求算卡或 CANN 安装，而是提供 starter code、环境清单、资料片段和阶段日志，使研究聚焦于协同使用知识与工程证据，而不是硬件性能。

本文贡献如下：
\begin{enumerate}[leftmargin=*]
  \item 提出 DACM，将 Agent 的上下文--检索--行动--反馈循环同开发者的目标设定、事实补充、判断与验证联系起来；
  \item 提出多主体协同轨迹分析，将开发过程表示为人、Agent、知识来源与本地环境间的有序事件，而非一个最终答案或满意度分数；
  \item 给出以形成性知识生态审计支撑质性协同研究的设计，并明确审计不能直接推断开发者行为；
  \item 提出面向开发者--Agent 协同体的技术社区设计启示。
\end{enumerate}

\section{相关研究}

\subsection{开发者如何与生成式 AI 协作}

生成式编程工具会改变开发者的起草、审阅、学习、信任与责任分配。开发者可能将 AI 建议当作草稿、对话伙伴、搜索捷径或快速试验工具，同时也会阅读 diff、运行测试、对照外部资料并限制权限来管理风险 \cite{barke2023grounded,vaithilingam2022expectation}。这些研究说明“采纳”不是一个二元决定。本文进一步将 Agent 对外部技术知识的依赖，以及开发者对本地事实的补充，纳入协同机制。

\subsection{Agent 工作循环及其边界}

与单次代码补全不同，Coding Agent 可以交织推理、检索、文件读取、编辑与工具调用 \cite{yao2023react,yang2024sweagent}。这一循环可将静态知识转化为候选行动，但受制于可见上下文、可访问证据、权限和反馈。本文将它理解为可观察的互动结构：它不是对所有模型内部推理方式的主张，而是用于观察开发者何时需要补充事实、选择来源、接受建议或亲自接管的分析基础。

\subsection{开发者知识实践与技术社区}

文档、样例、搜索、问答站点与同行网络长期影响开发者的学习和排障。在 Agentic Coding 中，这些资源也成为 Agent 的证据来源与校验器。本文将官方文档、样例、版本信息、仓库与第三方讨论共同称为\emph{技术知识生态}，由此提出传统门户评价之外的问题：知识能否被开发者与 Agent 共同发现、取得、归因并行动化？

\section{DACM：开发者--Agent 协同机制}

DACM 包含四类参与者与一个反馈循环。开发者设定目标和约束，决定可共享和可执行范围，评估风险并保持责任；Agent 建构上下文、检索和比较来源、提出计划或代码、调用被授权的工具并解释反馈；技术知识生态提供官方文档、样例、版本信息、仓库和社区讨论；本地工程环境包含代码、依赖、硬件、配置、数据、终端输出与测试，它是约束和最终验证发生之处。

\begin{figure}[t]
  \centering
  \Description{一个四主体循环：开发者、Coding Agent、技术知识生态与本地工程环境。箭头表示证据、候选行动、本地事实与验证反馈。}
  \fbox{\parbox{0.88\linewidth}{\centering\vspace{20mm}\textbf{图 1 占位：DACM 四主体闭环}\\
  开发者 $\leftrightarrow$ Agent $\leftrightarrow$ 技术知识生态 $\leftrightarrow$ 本地工程环境\\
  证据、候选行动、本地事实与验证反馈形成可回退的循环。\vspace{20mm}}}
  \caption{DACM 闭环。正式版本将替换为带有可访问描述的矢量图。}
  \label{fig:dacm-cn}
\end{figure}

框架通过五个维度落地：目标与上下文、知识取得、候选行动、现场执行、验证与回灌。每一维度都同时包含 Agent 的典型行动和开发者仍不可忽略的作用。比如，Agent 可以阅读兼容性页面，但开发者必须判断其版本范围是否适用于当前项目；Agent 可以运行测试，但开发者仍可能要判断结果是否满足质量与业务标准。

DACM 的最小分析单位是\emph{协同事件}：时间、发起者、行动、输入证据、输出状态、主导权变化和下一步去向。仅有“Agent 建议升级插件”并不足以描述协同；还需要记录建议依据什么来源，开发者是否提供版本表，是否真正执行，反馈如何回灌，以及谁决定停止。

本文不以“Agent 做了多少工作”衡量协同有效性。更有效的协同应能以可追溯证据推进任务、明确尚缺的现场事实、减少无效重复，并形成与风险相称的验证闭环。后续分析将从推进性、证据充分性、交接清晰度与验证可达性四个过程性结果比较案例。

\section{形成性知识生态审计（Study 0）}

Study 0 的作用是构造参与者研究的条件，而不是得出 CANN/CUDA 的总体排名。审计覆盖安装、算子开发、训练、推理部署、性能优化、调试与迁移等 26 个开发结果目标。每一目标分别在 CANN 与 CUDA 中使用真实术语构造问题对，在保持用户结果可比的前提下，由一个具备网页检索与读取能力的 Agent 按统一协议记录搜索、来源切换、官方正文、版本证据、第三方支撑、检索成本和最终答案的可行动性。

审计使用 11 项指标描述官方可发现性与可读取性、正文详尽度、版本清晰度、第三方支撑、模型先验、检索成本与答案产出。它与开发者门户旅程审计保持分离：前者描述 Agent 取得任务知识的条件，后者描述人类开发者直接进入生态时的基础体验。

审计产出三类形成性资产：任务条件卡、AddCustom 到 PyTorch 的连续工程情境包，以及两种知识供给构型。一类资料链相对连续、可读且版本明确；另一类需要在多份资料间拼接，或带有版本、正文可访问性和现场事实缺口。它们并非为了让参与者失败，而是为了观察在不能直接上卡时，开发者如何与 Agent 识别未知并设计验证。

\section{研究方法：开发者协同轨迹研究}

\subsection{参与者与流程}

计划招募约 20 名在近三个月实际使用过 AI Coding 工具的开发者。样本追求角色、经验、IDE/CLI 习惯、Agent 类型和任务领域的差异，而非统计代表性。参与者不需要拥有Ascend设备、维护 CANN 环境或提交完整私有仓库；他们可遮盖、跳过或删除任何材料。

单次研究包括：知情同意与背景问卷（10 分钟）；近期 AI 协同任务回放（35--45 分钟）；统一连续工程情境任务（50--60 分钟）；任务后回放与简短反思（15--20 分钟）；引用范围确认（5 分钟）。Study 1 可使用 Agent 对话、终端片段、代码 diff、Issue、PR、文档访问、截图或参与者叙述作为回放锚点。主持人沿时间追问已发生的行动，而不是要求参与者完成一份抽象计划作业。

Study 2 中，参与者在自己的工具中使用轻量情境包。五个连续节点是：D 写 Kernel、M Tiling、X 编译/内存证据、O 注册与框架接入、L 精度验收。starter code、环境清单、资料链接/片段和阶段日志构成可信工程上下文，但不假装提供真实算卡执行。参与者可以自然地搜索、询问常用 Agent、修改代码或记录、运行安全的本地脚本、说明不确定性，以及表达他们在真实环境中会验证什么。研究者仅管理时间和记录边界。

\subsection{数据与分析}

在取得同意后，研究收集访谈音频、必要的屏幕记录、Agent 对话、终端命令及输出、来源 URL、代码 diff、观察笔记与任务后回放。数据将去标识化；原始录音、转录、事件表和可引用片段分别设置访问权限。

分析结合多主体协同轨迹分析与反思性主题分析。事件字段包括发起者、动作、证据、知识状态、主导方、交接对象和结果。初始码本包含目标设定、上下文补充、官方/第三方/项目内知识、生成、检索、编辑、执行、事实请求、来源切换、人工接管、验证、回退和停止。两名研究者将先共同编码部分材料、讨论和修订定义，再持续比较跨案例模式。主题必须同时获得参与者叙述、轨迹事件和可用工件支持；Study 0 的评分不用于推断开发者行为。

\section{发现（待参与者研究完成后写入）}

\todo{本节在真实参与者数据完成前不应填入假设性结论。}

正式发现将围绕四个问题组织：（1）Agent 何时将找资料压缩为候选行动、何时反而增加核对成本；（2）开发者如何识别并提供本地事实；（3）知识供给条件如何改变协同轨迹形状，而不仅是一轮答案；（4）开发者如何描述自己在边界设定、证据判断与责任承担中的持续角色。每项发现至少包含跨案例模式、匿名引文、协同轨迹片段，以及一个反例或边界条件。

\section{设计启示}

\subsection{发布可行动的证据单元}
技术文档应围绕任务发布目标、适用版本、前置条件、最小命令或代码、预期输出、常见失败、恢复路径和稳定来源。这样的单元使 Agent 能引用而不是拼凑不受支持的步骤，也让开发者能快速判断建议是否适用。

\subsection{把版本关系变成可查询约束}
驱动、固件、工具链、框架、插件和硬件的配套关系往往是协同最早需要的现场事实。社区应提供机器可读、可深链的兼容矩阵和最小环境 manifest；Agent 在信息缺失时应明确请求事实，而非静默猜测。

\subsection{支持安全的事实回灌与验证}
开发者不一定能暴露完整环境，但可安全提供脱敏的版本摘要、日志片段、依赖树、测试结果和错误码。工具应解释“为什么需要此事实”“它影响哪个判断”。对于版本敏感或高风险修改，系统应在自动执行前呈现证据与最小验证 checkpoint。

\section{讨论与局限}

技术社区现在面对的是复合使用者：开发者与 Agent 共同进入、阅读、组织和验证技术知识。开发者仍是责任主体，但其劳动可能更集中于目标设定、事实补充、风险判断、验证与收束；Agent 则承担大量检索、压缩、起草和局部行动。这个分工是否有效，取决于知识能否以可靠证据的方式流动。

本研究有明确边界。Study 0 固定一个 Agent 与工具条件，不能代表全部模型或企业检索系统；CANN/CUDA 审计是形成性材料，不是硬件或生态的永久排名；约 20 人的质性研究适合机制发现，不适合总体比例或因果估计；统一情境包不能替代长期、真实上卡和团队协作中的观察。这些边界也指向后续的现场纵向研究、跨 Agent 比较和针对文档/工具设计的干预实验。

\section{结论}

DACM 将 Agentic Coding 描述为开发者、Agent、技术知识生态和本地工程环境之间的协同。通过追踪目标、证据、行动、现场事实、交接和验证，本文将“AI 是否有用”从单一答案质量问题，转化为可观察、可设计和可改进的人机协同机制问题。后续参与者研究将说明技术社区如何不只服务于人类读者，也服务于开发者--Agent 协同体。

\bibliographystyle{ACM-Reference-Format}
\bibliography{references}
\end{document}
